% arara: clean: {
% arara: --> extensions:
% arara: --> ['aux', 'log', 'nav', 'out',
% arara: --> 'pdf', 'snm', 'toc']
% arara: --> }
% arara: lualatex: {
% arara: --> shell: yes,
% arara: --> draft: yes,
% arara: --> interaction: batchmode
% arara: --> }
% arara: biber
% arara: lualatex: {
% arara: --> shell: yes,
% arara: --> draft: no,
% arara: --> interaction: batchmode
% arara: --> }
% arara: lualatex: {
% arara: --> shell: yes,
% arara: --> draft: no,
% arara: --> interaction: batchmode
% arara: --> }
% arara: clean: {
% arara: --> extensions:
% arara: --> ['aux', 'log', 'nav', 'out',
% arara: --> 'snm', 'toc']
% arara: --> }
\documentclass[
    12pt,
    aspectratio=1610,
    c,
    spanish,
    intlimits,
    handout,
    leqno,
    % noamsthm,
    professionalfonts,
]{beamer}
\usepackage[spanish]{babel}
\usetheme[
	style=academic,
	palette=oxford,
	titlealign=center,
	mathface=pagella,
	allserif]{Celestia}
\usepackage[
	citestyle=numeric,
	style=numeric,
	backend=biber,
]{biblatex}
\addbibresource{references.bib}

\setbeamertemplate{bibliography item}{%
	\ifboolexpr{ test {\ifentrytype{book}} or test {\ifentrytype{mvbook}}
		or test {\ifentrytype{collection}} or test {\ifentrytype{mvcollection}}
		or test {\ifentrytype{reference}} or test {\ifentrytype{mvreference}} }
	{\setbeamertemplate{bibliography item}[book]}
	{\ifentrytype{online}
		{\setbeamertemplate{bibliography item}[online]}
		{\setbeamertemplate{bibliography item}[article]}}%
	\usebeamertemplate{bibliography item}}

\defbibenvironment{bibliography}
{\list{}
	{\settowidth{\labelwidth}{\usebeamertemplate{bibliography item}}%
		\setlength{\leftmargin}{\labelwidth}%
		\setlength{\labelsep}{\biblabelsep}%
		\addtolength{\leftmargin}{\labelsep}%
		\setlength{\itemsep}{\bibitemsep}%
		\setlength{\parsep}{\bibparsep}}}
{\endlist}
{\item}

\title{Disminución de la dimensionalidad a través de proyecciones aleatorias}
\subtitle{Machine Learning para Matemáticos}
\author{Fidel Jara Huanca}
\institute{Instituto de Matemática y Ciencias Afines}
\date{23 de abril de 2026}

\begin{document}

\begin{frame}[plain]
	\titlepage
\end{frame}

\begin{frame}[label=toc]\transblindsvertical
	\frametitle{\contentsname}

	\tableofcontents[hideallsubsections]
\end{frame}

\section{Planteamiento del problema}

\begin{frame}
	\frametitle{\insertsection}

	\begin{alertblock}{}
		Describir, explicar el teorema de Johnson-Lindenstrauss que se
		justifica en reducción de la dimensionalidad en el contexto de
		machine learning.
		\begin{figure}[ht!]
			\centering
			\includegraphics[width=.65\paperwidth]{jl}
		\end{figure}
	\end{alertblock}
\end{frame}

\section{Revisión bibliográfica}

\begin{frame}
	\frametitle{\insertsection}
	Realizaremos un estudio acerca de
	\begin{itemize}
		\item

		      Maldición de la dimensionalidad.

		\item

		      Distribución uniforme en la bola unidad en alta dimensión.

		\item

		      Distribución normal multivariante en alta dimensión.
	\end{itemize}
	usando el texto de curso~\cite{Wegner2024}.
	\begin{refsection}
		\nocite{Wegner2024}
		\nocite{Johnson1984}
		\nocite{Dasgupta2003}
		\printbibliography[heading=none]
	\end{refsection}
\end{frame}

\section{Desarrollo matemático}

\begin{frame}
	\frametitle{\insertsection}

	\begin{itemize}
		\item

		      Teorema de la proyección aleatoria.

		\item

		      Teorema de Johnson-Lindenstrauss.
	\end{itemize}

\end{frame}

\section{Experimentos numéricos en Python}

\begin{frame}
	\frametitle{\insertsection}

	\begin{itemize}
		\item

		      Distorsión de las distancias proyectadas.

		\item

		      Comprobación del teorema Johnson-Lindenstrauss para distintos $\varepsilon$ y $k$.

		\item

		      Efecto de la dimensionalidad en $m$-vecinos más próximos.

		\item

		      Efecto de la dimensionalidad en $m$-medias.

		\item

		      Efectos de la dimensionalidad en análisis discriminatorio lineal (LDA) para separación y estimación paramétrica de distribuciones normales.

	\end{itemize}
\end{frame}

\end{document}
